\subsection*{直线和双曲线联立强化题}
\begin{enumerate}
    \item (2026$\cdot$上海春季卷$\cdot$20)已知双曲线 $\Gamma  : \dfrac{{x}^{2}}{2} - \dfrac{{y}^{2}}{2} = 1$ ,过点 $M\left( {m,0}\right)$ 作不垂直于 $x$ 轴的直线 $l$ 交双曲线 $\Gamma$ 于 $A$ 、 $B$ 两点.
    \begin{enumerate}
        \item 求双曲线离心率
        \item 若点 $A\left( {\sqrt{3},1}\right)$ ,点 $B$ 在双曲线的右支上,且 $B$ 是 ${AM}$ 的中点,求直线 $l$ 的斜率
        \item 若 $m > 0,{F}_{1},{F}_{2}$ 分别是双曲线 $\Gamma$ 的左右焦点, ${A}^{\prime }$ 是 $A$ 关于 $y$ 轴的对称点,若存在直线 $l$使得 $\vv{{F}_{1}{A}^{\prime }} \cdot  \vv{{F}_{2}B} = 0$ ,求 $m$ 的取值范围.
    \end{enumerate}
    \begin{jieda}
        \begin{enumerate}
            \item $c = 2, a = \sqrt{2}, e = \dfrac{c}{a} = \sqrt{2}$
            \item $B\left(\dfrac{\sqrt{3}+m}{2}, \dfrac{1}{2}\right)$在双曲线上，$\left(\dfrac{\sqrt{3}+m}{2}\right)^2 - \dfrac{1}{4} = 2$，故$(\sqrt{3}+m)^2 =9$又因为$B$在右支上，所以$\dfrac{\sqrt{3}+m}{2} > 0$，因此$m = 3-\sqrt{3}$，即$B\left(\dfrac{3}{2}, \dfrac{1}{2}\right)$，因此$k_l = \dfrac{1-\dfrac{1}{2}}{\sqrt{3}-\dfrac{3}{2}} = \dfrac{1}{2\sqrt{3}-3} = \dfrac{2\sqrt{3}+3}{3}$
            \item $F_1(-2, 0)$，$F_2(2, 0)$
            
            当$l$为$y = 0$时不满足要求，故设直线$l:  x = ty + m$，设$A(x_1, y_1), B(x_2, y_2)$，因此$A'(-x_1, y_1)$

            $\vv{F_1A'} = (2-x_1, y_1)$，$\vv{F_2A} = (x_2-2, y_2)$，$\vv{F_1A'} \cdot \vv{F_2A} =(2-x_1)(x_2-2)+y_1y_2 = 0$，即$(2-m-ty_1)(ty_2+m-2)+y_1y_2 = 0$，即$(1-t^2)y_1y_2+(2-m)t(y_1+y_2)-(m-2)^2 = 0$

            联立直线与双曲线$\left\{\begin{array}{l}
                 x = ty+m  \\
                 x^2-y^2 = 2
            \end{array}\right.$，消去$x$得$(ty+m)^2-y^2 = 2$，即$(t^2-1)y^2+2tmy+(m^2-2) = 0$，$t \neq \pm 1$，$\Delta = 4t^2m^2 - 4(m^2-2)(t^2-1) = 4m^2+8t^2-8 > 0$

            根据韦达定理有$\left\{\begin{array}{l}
                 y_1+y_2 = \dfrac{2tm}{1-t^2} \\
                 y_1y_2 = \dfrac{2-m^2}{1-t^2}
            \end{array}\right.$，代入$(1-t^2)y_1y_2+(2-m)t(y_1+y_2) -(m-2)^2 = 0$得$(1-t^2)\cdot\dfrac{2-m^2}{1-t^2}+(2-m)t\cdot\dfrac{2tm}{1-t^2}-(m-2)^2 = 0$，即$\dfrac{t^2-m^2+2m-1}{1-t^2} = 0$，故$t^2 = (m-1)^2$

            代入$4m^2+8t^2-8 > 0$可得$m^2+2(m-1)^2-2>0$，解得$m < 0$或$m > \dfrac{4}{3}$，又因为$m > 0, t^2 \neq 1$，所以$m \in \left(\dfrac{4}{3}, 2\right) \cup (2, +\infty)$
        \end{enumerate}
    \end{jieda}
\end{enumerate}
\subsection*{抛物线定点与韦达结论、面积之比转化}
\begin{enumerate}
    \item (2024$\cdot$闵行区一模$\cdot$20)如图,已知椭圆 ${C}_{1} : \dfrac{{x}^{2}}{4} + {y}^{2} = 1$ 和抛物线 ${C}_{2} : {x}^{2} = {2py}\left( {p > 0}\right) ,{C}_{2}$ 的焦点 $F$ 是 ${C}_{1}$ 的上顶点，过 $F$ 的直线交 ${C}_{2}$ 于 $M$ 、 $N$ 两点，连接 ${NO}$ 、 ${MO}$ 并延长之，分别交 ${C}_{1}$ 于 $A$ 、 $B$ 两点， 连接 ${AB}$ ，设 $\bigtriangleup  {OMN}$ 、 $\bigtriangleup  {OAB}$ 的面积分别为 ${S}_{\bigtriangleup {OMN}}$ 、 ${S}_{\bigtriangleup {OAB}}$ .
    \begin{center}  
    \begin{minipage}{0.45\textwidth}
            \begin{enumerate}
        \item 求 $p$ 的值
        \item 求 $\vv{OM} \cdot  \vv{ON}$ 的值
        \item 求 $\dfrac{{S}_{\bigtriangleup {OMN}}}{{S}_{\bigtriangleup {OAB}}}$ 的取值范围.
    \end{enumerate}
    \end{minipage}
    \begin{minipage}{0.45\textwidth}
        \includegraphics[width=0.45\textwidth]{bodyxb/src/0_238_1291_365_283_0.jpg}
    \end{minipage}  
    \end{center}
    \begin{jieda}
        \begin{enumerate}
            \item $F(0, 1)$，$p= 2$
            \item 设$M(x_1, y_1), N(x_2, y_2)$，则$\vv{OM} \cdot \vv{ON} = x_1x_2+y_1y_2 = x_1x_2 + \dfrac{x_1^2x_2^2}{16}$
            
            易知$MN$斜率存在，故设$MN: y = kx+1$，与抛物线联立可得$x^2 - 4kx-4 = 0$，$\Delta > 0$，根据韦达定理可知$x_1x_2 = -4$，故$\vv{OM} \cdot \vv{ON} = -3$
            \item $\dfrac{{S}_{\bigtriangleup {OMN}}}{{S}_{\bigtriangleup {OAB}}} = \dfrac{|OM||ON|}{|OA||OB|} = \dfrac{\vv{OM} \cdot \vv{ON}}{\vv{OA} \cdot \vv{OB}} = \dfrac{-3}{\vv{OA} \cdot \vv{OB}}$

            设直线 ${NO}$ 、 ${MO}$ 的斜率分别为 ${k}_{1}$ 、 ${k}_{2}$ ,不妨设 ${k}_{1} > 0,{k}_{2} < 0$，$\vv{OA} \cdot \vv{OB} = x_Ax_B(1+k_1k_2)$

            $k_1k_2 = \dfrac{y_1y_2}{x_1x_2} = \dfrac{x_1x_2}{16} = -\dfrac{1}{4}$
            
            
            联立 $\left\{  \begin{array}{l} y = {k}_{1}x \\  {x}^{2} + 4{y}^{2} = 4 \end{array}\right.$ 可得 $\left( {4{k}_{1}^{2} + 1}\right) {x}^{2} = 4$ ,解得 $x =  \pm  \dfrac{2}{\sqrt{4{k}_{1}^{2} + 1}}$ 
            
            点 $A$ 在第三象限,则 ${x}_{A} =  - \dfrac{2}{\sqrt{4{k}_{1}^{2} + 1}}$，点 $B$ 在第四象限,同理可得 ${x}_{B} = \dfrac{2}{\sqrt{4{k}_{2}^{2} + 1}} = \dfrac{2}{\sqrt{\dfrac{4}{16k_1^2}+1}} = \dfrac{4k_1}{\sqrt{1+4k_1^2}}$ 

            故$\vv{OA} \cdot \vv{OB} = \dfrac{-6k_1}{1+4k_1^2}$，故$\dfrac{{S}_{\bigtriangleup {OMN}}}{{S}_{\bigtriangleup {OAB}}} = \dfrac{1+4k_1^2}{2k_1} = 2k_1+\dfrac{1}{2k_1} \in [2,+\infty)$
            
        \end{enumerate}
    \end{jieda}
\end{enumerate}
\subsection*{分式型非对称韦达}
\begin{enumerate}
    \item (2025$\cdot$同济大学一附中高三9月月考$\cdot$20(3))已知在平面直角坐标系 ${xOy}$ 中,椭圆 $C : \dfrac{{x}^{2}}{4} + \dfrac{{y}^{2}}{3} = 1$ ,其右焦点为 $F$ ,过点 $F$ 且与坐标轴不垂直的直线 $l$ 与椭圆 $C$ 交于 $P,Q$ 两点，若 $\vv{OF} = \dfrac{2}{5}\vv{OP} + \dfrac{3}{5}\vv{OQ}$ ，求直线 $l$ 的方程.
    \begin{jieda}
        \begin{enumerate}
            \item 方法1(极菜)
                
            由题可知，直线 $l$ 的斜率存在且不为零，设直线 $l$ 的方程为 $y = k\left( {x - 1}\right)$ ，
            
            设 $P\left( {{x}_{1},k\left( {{x}_{1} - 1}\right) }\right) ,Q\left( {{x}_{2},k\left( {{x}_{2} - 1}\right) }\right)$ ,
            
            因为 $\vv{OF} = \dfrac{2}{5}\vv{OP} + \dfrac{3}{5}\vv{OQ}$ ,所以 $\left( {1,0}\right)  = \dfrac{2}{5}\left( {{x}_{1},k\left( {{x}_{1} - 1}\right) }\right)  + \dfrac{3}{5}\left( {{x}_{2},k\left( {{x}_{2} - 1}\right) }\right)$ .
            
            所以有 $\left\{  \begin{array}{l} \dfrac{2}{5}{x}_{1} + \dfrac{3}{5}{x}_{2} = 1 \\  \dfrac{2}{5}k\left( {{x}_{1} - 1}\right)  + \dfrac{3}{5}k\left( {{x}_{2} - 1}\right)  = 0 \end{array}\right.$ ,化简得 $2{x}_{1} + 3{x}_{2} = 5$ .

            联立直线方程与椭圆方程得 $3{x}^{2} + 4{k}^{2}{\left( x - 1\right) }^{2} = {12}$ ,
            
            化简得 $\left( {4{k}^{2} + 3}\right) {x}^{2} - 8{k}^{2}x + 4{k}^{2} - {12} = 0$ .
            
            $\Delta  = {64}{k}^{4} - 4\left( {4{k}^{2} + 3}\right) \left( {4{k}^{2} - {12}}\right)  = {144}{k}^{2} + {144},$
            
            所以解得 $x = \dfrac{8{k}^{2} \pm  {12}\sqrt{{k}^{2} + 1}}{2\left( {4{k}^{2} + 3}\right) } = \dfrac{4{k}^{2} \pm  6\sqrt{{k}^{2} + 1}}{4{k}^{2} + 3}$
            
            若 ${x}_{1} = \dfrac{4{k}^{2} + 6\sqrt{{k}^{2} + 1}}{4{k}^{2} + 3},{x}_{2} = \dfrac{4{k}^{2} - 6\sqrt{{k}^{2} + 1}}{4{k}^{2} + 3}$，将其代入 $2{x}_{1} + 3{x}_{2} = 5$ 中得 $2 \times  \dfrac{4{k}^{2} + 6\sqrt{{k}^{2} + 1}}{4{k}^{2} + 3} + 3 \times  \dfrac{4{k}^{2} - 6\sqrt{{k}^{2} + 1}}{4{k}^{2} + 3} = 5$ ,化简得 $- 2\sqrt{{k}^{2} + 1} = 5$ ,无解.
            
            若 ${x}_{1} = \dfrac{4{k}^{2} - 6\sqrt{{k}^{2} + 1}}{4{k}^{2} + 3},{x}_{2} = \dfrac{4{k}^{2} + 6\sqrt{{k}^{2} + 1}}{4{k}^{2} + 3}$ ,
            
            将其代入 $2{x}_{1} + 3{x}_{2} = 5$ 中得$2 \times  \dfrac{4{k}^{2} - 6\sqrt{{k}^{2} + 1}}{4{k}^{2} + 3} + 3 \times  \dfrac{4{k}^{2} + 6\sqrt{{k}^{2} + 1}}{4{k}^{2} + 3} = 5$ ,化简得 $2\sqrt{{k}^{2} + 1} = 5$
            
            解得 $k =  \pm  \dfrac{\sqrt{21}}{2}$ ,所以直线方程为 $y =  \pm  \dfrac{\sqrt{21}}{2}\left( {x - 1}\right)$ .

            \item 方法2

            $\vv{OF} = \dfrac{2}{5}\vv{OP} + \dfrac{3}{5}\vv{OQ} \Longleftrightarrow \dfrac{|FP|}{|FQ|} = \dfrac{3}{2}$

            设$PQ$的倾斜角为$\alpha$，利用椭圆的焦半径公式
            \begin{dingautolist}{172}
                \item 若$\left\{\begin{array}{l}
                        |FP| = \dfrac{ep}{1-ecos\alpha}  \\
                        |FQ| = \dfrac{ep}{1+ecos\alpha}
                \end{array}\right.$，则$\dfrac{1+ecos\alpha}{1-ecos\alpha} = \dfrac{3}{2}$，故$2+2ecos\alpha = 3-3ecos\alpha$，故$cos\alpha = \dfrac{1}{5e} = \dfrac{2}{5}$
    
                \item  若$\left\{\begin{array}{l}
                        |FP| = \dfrac{ep}{1+ecos\alpha}  \\
                        |FQ| = \dfrac{ep}{1-ecos\alpha}
                \end{array}\right.$，则$\dfrac{1-ecos\alpha}{1+ecos\alpha} = \dfrac{3}{2}$，故$3+3ecos\alpha = 2-2ecos\alpha$，故$cos\alpha = -\dfrac{1}{5e} = -\dfrac{2}{5}$
            \end{dingautolist}
            故$k = \pm \dfrac{\sqrt{21}}{2}$,所以直线方程为 $y =  \pm  \dfrac{\sqrt{21}}{2}\left( {x - 1}\right)$ .

            \item 方法3
            
            $\vv{OF} = \dfrac{2}{5}\vv{OP} + \dfrac{3}{5}\vv{OQ} \Longleftrightarrow \dfrac{|FP|}{|FQ|} = \dfrac{3}{2} \Longleftrightarrow \dfrac{y_P}{y_Q} = -\dfrac{3}{2}$，故$\dfrac{y_Q}{y_P} = -\dfrac{2}{3}$，故$\dfrac{y_P}{y_Q} + \dfrac{y_Q}{y_P} = \dfrac{y_P^2+y_Q^2}{y_Py_Q} = \dfrac{(y_P+y_Q)^2-2y_Py_Q}{y_Py_Q} = \dfrac{(y_P+y_Q)^2}{y_Py_Q} - 2 = -\dfrac{13}{6}$

            易知直线斜率为0时不符合要求，设$PQ: x = my + 1$，联立直线和椭圆方程$\left\{\begin{array}{l}
                    \dfrac{{x}^{2}}{4} + \dfrac{{y}^{2}}{3} = 1  \\
                    x = my + 1
            \end{array}\right.$，则$3(my + 1)^2 + 4y^2 = 12$，即$(3m^2+4)y^2+6my-9=0$，故$\left\{\begin{array}{l}
                    y_P+y_Q = \dfrac{-6m}{3m^2+4} \\
                    y_Py_Q = \dfrac{-9}{3m^2+4}
            \end{array}\right.$

            $\Delta > 0$，故有$\dfrac{(y_P+y_Q)^2}{y_Py_Q} - 2 = \dfrac{\left(\dfrac{-6m}{3m^2+4}\right)^2}{\dfrac{-9}{3m^2+4}} - 2 = \dfrac{-4m^2}{3m^2+4} - 2 = -\dfrac{13}{6}$

            即$\dfrac{-4m^2}{3m^2+4} = -\dfrac{1}{6}$，即$3m^2+4 = 24m^2$，故$m = \pm \dfrac{2}{\sqrt{21}}$, 所以直线方程为 $y =  \pm  \dfrac{\sqrt{21}}{2}\left( {x - 1}\right)$ .
        \end{enumerate}
    \end{jieda}
\end{enumerate}
\subsection*{斜率之和、斜率之积为定值}
\begin{enumerate}
    \item 已知双曲线 $C : \dfrac{{x}^{2}}{{a}^{2}} - \dfrac{{y}^{2}}{{b}^{2}} = 1\left( {a > 0,b > 0}\right)$ 的离心率为 $\dfrac{\sqrt{6}}{2}$ ,右顶点为 $E\left( {\sqrt{2},0}\right) .A,B$ 为双曲线 $C$ 右支上两点,且点 $A$ 在第一象限,以 ${AB}$ 为直径的圆经过点 $E$.
    \begin{center}  
    \begin{minipage}{0.45\textwidth}
    \begin{enumerate}
        \item 求 $C$ 的方程
        \item 证明:直线 ${AB}$ 恒过定点
        \item 若直线 ${AB}$ 与 $x,y$ 轴分别交于点 $M,P$且 $M$ 为 ${PA}$ 中点,求 $\dfrac{{S}_{\bigtriangleup {PBE}}}{{S}_{\bigtriangleup {MBE}}}$ 的值
    \end{enumerate}
    \end{minipage}
    \begin{minipage}{0.45\textwidth}
        \includegraphics[width=0.45\textwidth]{bodyxb/src/0_1147_1875_368_354_0.jpg}
    \end{minipage}  
    \end{center}
    \begin{jieda}
        \begin{enumerate}
            \item $\dfrac{x^2}{2} - y^2 = 1$
            \item 设 $A\left( {{x}_{1},{y}_{1}}\right) ,B\left( {{x}_{2},{y}_{2}}\right)$ ,可设直线 ${AB} : x = {my} + t$ . 联立 $\left\{  \begin{array}{l} x = {my} + t \\  \dfrac{{x}^{2}}{2} - {y}^{2} = 1 \end{array}\right.$ ,得 $\left( {{m}^{2} - 2}\right) {y}^{2} + {2mty} + {t}^{2} -2 = 0$ ,则 $\left\{  \begin{matrix} {m}^{2} - 2 \neq  0 \\  \Delta  =  - 8\left( {{m}^{2} + {t}^{2} - 2}\right)  > 0 \end{matrix}\right.$ ,即 $\left\{ \begin{matrix} {m}^{2} \neq  2 \\  {m}^{2} + {t}^{2} < 2 \end{matrix}\right.$
            
            根据韦达定理有$\left\{\begin{array}{l}
                {y}_{1} + {y}_{2} =  - \dfrac{2mt}{{m}^{2} - 2} \\
                {y}_{1}{y}_{2} = \dfrac{{t}^{2} - 2}{{m}^{2} - 2}
            \end{array}\right.$.
            
            $\vv{AE} \cdot  \vv{BE} = 0$，即$\left( {{m}^{2} + 1}\right) {y}_{1}{y}_{2} + m\left( {t - \sqrt{2}}\right) \left( {{y}_{1} + {y}_{2}}\right)  + {\left( t - \sqrt{2}\right) }^{2} = 0$

            代入得$\dfrac{\left( {{m}^{2} + 1}\right) \left( {{t}^{2} - 2}\right) }{{m}^{2} - 2} - \dfrac{m\left( {t - \sqrt{2}}\right)  \cdot  {2mt}}{{m}^{2} - 2} + {\left( t - \sqrt{2}\right) }^{2} = 0$，即$\left( {t - \sqrt{2}}\right) \left( {3\sqrt{2} - t}\right)  = 0$
            
            当 $t = \sqrt{2}$ 时,直线 ${AB} : x = {my} + \sqrt{2}$ 经过点 $E$ ,不符条件,舍去.
            
            故直线 ${AB} : x = {my} + 3\sqrt{2}$ 必过定点 $M\left( {3\sqrt{2},0}\right)$ .
            
            \item $P\left( {0, - \dfrac{3\sqrt{2}}{m}}\right) ,M\left( {3\sqrt{2},0}\right) ,M$ 为 ${PA}$ 中点, $\therefore A\left( {6\sqrt{2},\dfrac{3\sqrt{2}}{m}}\right)$ ,代入双曲线得 ${m}^{2} = \dfrac{18}{35}$.
            
            由 ${y}_{1}{y}_{2} = \dfrac{3\sqrt{2}}{m}{y}_{2} = \dfrac{16}{{m}^{2} - 2}$ 得 ${y}_{2} = \dfrac{8\sqrt{2}m}{3\left( {{m}^{2} - 2}\right) }$.

            $\dfrac{{S}_{\bigtriangleup {PBE}} + {S}_{\bigtriangleup {MBE}}}{{S}_{\bigtriangleup {MBE}}} = \dfrac{{S}_{\bigtriangleup {PEM}}}{{S}_{\bigtriangleup {MBE}}} = \dfrac{\left| {y}_{P}\right| }{\left| {y}_{2}\right| } = \dfrac{\left| \dfrac{3\sqrt{2}}{m}\right| }{\left| {y}_{2}\right| } = \left| \dfrac{9\left( {{m}^{2} - 2}\right) }{8{m}^{2}}\right|  = \dfrac{13}{4},\therefore \dfrac{{S}_{\bigtriangleup {BEP}}}{{S}_{\bigtriangleup {BEM}}} = \dfrac{9}{4}$
        \end{enumerate}
    \end{jieda}

    \item （选做）(2024$\cdot$闵行区高二下期末$\cdot$20)已知椭圆 $\Gamma  : \dfrac{{x}^{2}}{2} + {y}^{2} = 1$ 的左、右焦点分别为 ${F}_{1},{F}_{2},A,B$ 分别为 $\Gamma$ 的上,下顶点,$P\left( {{x}_{1},{y}_{1}}\right) ,Q\left( {{x}_{2},{y}_{2}}\right)$ 是 $\Gamma$ 上不同于点 $A$ 的两点.
    \begin{enumerate}
        \item 求 $\left| {{F}_{1}{F}_{2}}\right|$ 的值
        \item 记 $\bigtriangleup P{F}_{1}{F}_{2},\bigtriangleup {PAB}$ 的面积分别为 ${S}_{1},{S}_{2}$ ,若 ${S}_{1} < {S}_{2}$ ,求 $x_{1}$ 的取值范围
        \item 若直线 ${AP}$ 与 ${AQ}$ 的斜率之和为 2,作 ${AH} \bot  {PQ}$，垂足为 $H$ ,试问: 点 $H$ 是否在一个定圆上? 若是,求出该圆的方程; 若不是,说明理由.
    \end{enumerate}
    \begin{jieda}
        \begin{enumerate}
            \item $2c = 2$
            \item $S_1 = |y_1|$，$S_2 = |x_1|$，则$S_1 < S_2$即$|y_1| < |x_1|$，即$y_1^2 < x_1^2$，故$1-\dfrac{x_1^2}{2} < x_1^2$，因此$x_1 \in \left(-\sqrt{2}, -\dfrac{\sqrt{6}}{3}\right) \cup \left(\dfrac{\sqrt{6}}{3}, \sqrt{2}\right)$
            \item 设$PQ: y = kx+m (m \neq 1)$，联立方程 $\left\{  \begin{array}{l} y = {kx} + m \\  \dfrac{{x}^{2}}{2} + {y}^{2} = 1 \end{array}\right.$ ,消去 $y$ 得 $\left( {2{k}^{2} + 1}\right) {x}^{2} + {4kmx} + 2{m}^{2} - 2 = 0$，则 $\Delta  > 0$ ,可得 $\left\{\begin{array}{l}
                {x}_{1} + {x}_{2} =  - \dfrac{4km}{2{k}^{2} + 1}  \\
                {x}_{1}{x}_{2} = \dfrac{2{m}^{2} - 2}{2{k}^{2} + 1}
            \end{array}\right.$

            ${k}_{AP} + {k}_{AQ} = \dfrac{{y}_{1} - 1}{{x}_{1}} + \dfrac{{y}_{2} - 1}{{x}_{2}} = \dfrac{k{x}_{1} + m - 1}{{x}_{1}} + \dfrac{k{x}_{2} + m - 1}{{x}_{2}} = 2k + (m-1)\left(\dfrac{x_1+x_2}{x_1x_2}\right) = 2k + (m-1)\left(\dfrac{-4km}{2{m}^{2} - 2}\right) = 2k - \dfrac{2km}{m+1} = \dfrac{2k}{m+1} = 2$，故$k = m+1$，即$PQ$过定点$M(-1, -1)$

            若直线 ${PQ}$ 的斜率不存在,则 $P\left( {{x}_{1},{y}_{1}}\right) ,Q\left( {{x}_{1}, - {y}_{1}}\right)$，因为直线 ${AP}$ 与 ${AQ}$ 的斜率之和为 2，则 ${k}_{AP} + {k}_{AQ} = \dfrac{{y}_{1} - 1}{{x}_{1}} + \dfrac{-{y}_{1} - 1}{{x}_{1}} =  - \dfrac{2}{{x}_{1}} = 2$ ,即 ${x}_{1} =  - 1$，此时直线 ${PQ} : x =  - 1$ 也过定点 $M\left( {-1, - 1}\right)$，综上，直线 ${PQ}$ 过定点 $M\left( {-1, - 1}\right)$

            故$H$在以$AM$为直径的圆$(x-0)(x+1)+(y-1)(y+1) = 0$，即$x^2+x+y^2 = 1$上.
        \end{enumerate}
    \end{jieda}

    \item （选做）过抛物线 ${y^2} = {4x}$ 上点 $P\left( {1,2}\right)$ 作三条斜率分别为 ${k}_{1},{k}_{2},{k}_{3}$ 的直线 ${l}_{1},{l}_{2},{l}_{3}$ . 与抛物线分别交于不同于 $p$ 的点 $A,B,C$ . 若 ${k}_{1} + {k}_{2} = 0,{k}_{2} \cdot  {k}_{3} =  - 1$，则以下结论正确的是\dotfill{(\kh{B})}
    \xx{直线 ${AB}$ 过定点}{直线 ${AB}$ 斜率一定}{直线 ${BC}$ 斜率一定}{直线 ${AC}$ 斜率一定}
    \begin{jieda}
        由题意, ${k}_{1},{k}_{2},{k}_{3}$ 均不为 $0$,设$A\left( {{x}_{1},{y}_{1}}\right) ,B\left( {{x}_{2},{y}_{2}}\right) ,C\left( {{x}_{3},{y}_{3}}\right)$，则${k}_{1} = \dfrac{{y}_{1} - 2}{{x}_{1} - 1} = \dfrac{{y}_{1} - 2}{\dfrac{{y}_{1}^{2}}{4} - 1} = \dfrac{4}{{y}_{1} + 2}$ ,同理可得 ${k}_{2} = \dfrac{{y}_{2} - 2}{{x}_{2} - 1} = \dfrac{4}{{y}_{2} + 2}$ ,
        ${k}_{3} = \dfrac{{y}_{3} - 2}{{x}_{3} - 1} = \dfrac{4}{{y}_{3} + 2}$

        由 ${k}_{1} + {k}_{2} = 0$ ,得 $\dfrac{4}{{y}_{1} + 2} + \dfrac{4}{{y}_{2} + 2} = 0$ ,即 ${y}_{1} + {y}_{2} + 4 = 0$. 
        设直线 ${AB}$ 的方程为 $x = {m}_{1}y + {t}_{1}$ ,联立抛物线方程可得 ${y^2} - 4{m}_{1}y - 4{t}_{1} = 0$. 
        则 $\Delta  > 0, \left\{\begin{array}{l}
             {y}_{1} + {y}_{2} = 4{m}_{1}  \\
             {y}_{1}{y}_{2} =  - 4{t}_{1} 
        \end{array}\right.$ 代入${y}_{1} + {y}_{2} + 4 = 0$可得 $4{m}_{1} + 4 = 0,{m}_{1} =  - 1$，此时直线 ${AB}$ 的方程为 $x =  - y + {t}_{1}$ ,故直线 ${AB}$ 斜率是定值,故 $\mathrm{B}$ 正确, $\mathrm{A}$ 错误

        由 ${k}_{2} \cdot  {k}_{3} =  - 1$ ,得 $\dfrac{4}{{y}_{2} + 2} \times  \dfrac{4}{{y}_{3} + 2} =  - 1$ ,即 ${y}_{2}{y}_{3} + 2\left( {{y}_{2} + {y}_{3}}\right)  + {20} = 0$.
        设直线${BC}$ 的方程为 $x = {m}_{2}y + {t}_{2}$ ,联立抛物线方程可得 ${y}^{2} - 4{m}_{2}y - 4{t}_{2} = 0$.
        则 $\Delta  > 0, \left\{\begin{array}{l}
             {y}_{3} + {y}_{2} = 4{m}_{2}  \\
             {y}_{3}{y}_{2} =  - 4{t}_{2} 
        \end{array}\right.$，代入${y}_{2}{y}_{3} + 2\left( {{y}_{2} + {y}_{3}}\right)  + {20} = 0$可得 $2{m}_{2} - {t}_{2} + 5 = 0$ ,此时 ${BC}$ 的方程为$x = {m}_{2}y + 2{m}_{2} + 5 = {m}_{2}\left( {y + 2}\right)  + 5$ ,恒过定点$(5, -2)$,斜率不是定值,故C错误

        由 ${k}_{2} \cdot  {k}_{3} =  - 1,{k}_{1} + {k}_{2} = 0$ ,得 ${k}_{1} \cdot  {k}_{3} = 1$ ,即 $\dfrac{4}{{y}_{1} + 2} \times  \dfrac{4}{{y}_{3} + 2} = 1$，即 ${y}_{1}{y}_{3} + 2\left( {{y}_{1} + {y}_{3}}\right)  - {12} = 0$.
        设直线 ${AC}$ 的方程为 $x = {m}_{3}y + {t}_{3}$ ,联立抛物线方程可得 ${y}^{2} - 4{m}_{3}y - 4{t}_{3} = 0$ ,则 $\Delta  > 0, \left\{\begin{array}{l}
             {y}_{3} + {y}_{1} = 4{m}_{3}  \\
             {y}_{3}{y}_{1} =  - 4{t}_{3} 
        \end{array}\right.$代入${y}_{1}{y}_{3} + 2\left( {{y}_{1} + {y}_{3}}\right)  - {12} = 0$可得$2{m}_{3} - {t}_{3} - 3 = 0$ ,此时 ${AC}$ 的方程为 $x = {m}_{3}y + 2{m}_{3} - 3 = {m}_{2}\left( {y + 2}\right)  - 3$ 恒过定点$(-3, -2)$,斜率不为定值，故D错误.

        \ps 过抛物线$y^2 = 2px$上一点$P(a, b)$作两直线分别交抛物线于$A, B$，若两直线斜率之和为0，则$k_{AB} = -\dfrac{p}{b}$ \\
        \ps 过抛物线$y^2 = 2px$上一点$P(a, b)$作两直线分别交抛物线于$A, B$，若两直线斜率之积为常数$c$，则直线$AB$恒过定点$\left(a-\dfrac{2p}{c}, -b\right)$
    \end{jieda}
\end{enumerate}

\subsection*{端点效应}
\begin{enumerate}
    \item 若对任意$x \geq 1$，不等式$x\left(lnx-ax\right)+a \leq 0$恒成立，求实数$a$的取值范围.
        \begin{jieda}
            不等式化为$xlnx-ax^2+a\leq 0$
            \begin{enumerate}
                \item 方法1 分类讨论
                
                设$f(x) = xlnx-ax^2+a$，则$f'(x) = lnx+1-2ax$，$f''(x) = \dfrac{1}{x}-2a$

                \begin{dingautolist}{172}
                    \item 若$a \leq 0$，则$f''(x) > 0$恒成立，故$f'(x)$严增，故$f'(x) \geq f'(1) = 1-2a > 0$，故$f(x)$严增，故$f(x) \geq f(1)$，不符合要求
                    \item 若$a \geq \dfrac{1}{2}$，则$f''(x) \leq 0$恒成立，故$f'(x)$严减，故$f'(x) \leq f'(1) = 1-2a \leq 0$，故$f(x)$严减，故$f(x) \leq f(1)$，满足要求
                    \item 若$a \in \left(0, \dfrac{1}{2}\right)$，则$f''(x)$在$\left(1, \dfrac{1}{2a}\right)$大于0，故$f'(x)$在$\left(1, \dfrac{1}{2a}\right)$严增，故$f'(x)$在$\left(1, \dfrac{1}{2a}\right)$大于0，故$f(x)$在$\left(1, \dfrac{1}{2a}\right)$严增，故$f(x)$在$\left(1, \dfrac{1}{2a}\right)$大于0，不符合要求.
                \end{dingautolist}
                综上，$a \in \left[\dfrac{1}{2}, +\infty\right)$
            
                \item 方法2 参变分离

                当$x = 1$时恒成立，参变分离有$a \geq \dfrac{xlnx}{x^2-1}(x > 1)$，设$f(x) = \dfrac{xlnx}{x^2-1}$，则$f'(x) = \dfrac{(lnx+1)(x^2-1)-2x^2lnx}{(x^2-1)^2} = \dfrac{x^2-1-(1+x^2)lnx}{(x^2-1)^2}$

                令$g(x) = x^2-1-(1+x^2)lnx$，$g'(x) = 2x-2xlnx-\dfrac{1+x^2}{x} = x-2xlnx-\dfrac{1}{x}$，$g''(x) = 1+\dfrac{1}{x^2}-2(lnx+1)$，$g'''(x) = -\dfrac{2}{x^3}-\dfrac{2}{x} < 0$，故$g''(x) < 0$，故$g'(x) < 0$，故$g(x) < 0$，故$f'(x) < 0$，根据洛必达法则可知$f(x) < \dfrac{1}{2}$，故$a \geq \dfrac{1}{2}$
            \end{enumerate}
        \end{jieda}
\end{enumerate}